% No preamble. Prefix aqb-. Infinite analytic interactions; classical descent only.
\section{A Banach quotient for analytic interactions}
\label{aqb-section}

Fix the cell size $J$, a bound $D$ on the number of scalar field
coordinates per cell, a cardinality exponent $a>0$, and a radius
$b>0$. Let $\mathcal C$ be the countable set of nonempty finite
connected cell sets. A declared carrier may exceed the variables
actually read by its coefficient. Let $H^\infty_b(H)$ denote the
bounded holomorphic functions on the open scalar-coordinate
radius-$b$ polydisc of $H$. Define
\begin{equation}
 \mathcal A_{a,b}=\left\{E=(E_H)_{H\in\mathcal C}:
 \|E\|_{a,b}:=\sup_i\sum_{H\ni i}e^{a|H|}
                       \|E_H\|_b<\infty\right\}.
 \label{aqb-activity-space}
\end{equation}
Individual coefficients need not have any symmetry. Constants
are allowed. All constants below may depend on the fixed $J,D,a$.

\subsection{Closed equality kernel and summable fields}
\label{aqb-quotient-section}

For a finitely supported field $A$, $\|A\|_\infty<b$, put
\begin{equation}
 \Phi_E(A)=\sum_H\bigl[E_H(A|_H)-E_H(0)\bigr].
 \label{aqb-total}
\end{equation}
If $K$ is its finite cell support, only carriers meeting $K$
contribute, and
\begin{equation}
 \sum_H|E_H(A|_H)-E_H(0)|\le2|K|\|E\|_{a,b}.
 \label{aqb-evaluation-bound}
\end{equation}
The series is normally convergent on every finite-coordinate
polydisc.

\begin{proposition}[Complete representation quotient]
\label{aqb-quotient-proposition}
The space $\mathcal A_{a,b}$ is Banach. The subspace
\begin{equation}
 \mathcal N_{a,b}=\{E:\Phi_E(A)=0
                  \text{ for every compact }A,\ \|A\|_\infty<b\}
 \label{aqb-null-space}
\end{equation}
is closed, and
\begin{equation}
 \mathcal Q_{a,b}=\mathcal A_{a,b}/\mathcal N_{a,b},
 \qquad
 \|[E]\|_{\mathcal Q}=\inf_{F\in\mathcal N_{a,b}}
                                      \|E+F\|_{a,b}
 \label{aqb-quotient-norm}
\end{equation}
is a Banach space. Its invariant subspace, defined by translation,
hypercubic, global-adjoint and local-gauge invariance of the
normalized total, is closed.
\end{proposition}

\begin{proof}
A norm-Cauchy sequence is Cauchy in each $H^\infty_b(H)$:
choose one root in $H$ in \eqref{aqb-activity-space}. Component
limits exist. Fatou over carriers, followed by the root supremum,
proves both the norm bound and norm convergence. The null space
is the intersection of the kernels of the bounded evaluations
\eqref{aqb-evaluation-bound}. The usual quotient-completeness
argument uses a summably Cauchy subsequence and representatives
with summable successive differences; it requires no minimizing
representative. Each total symmetry is an equality of bounded
evaluations, for all compact configurations and transforms in
their common domain. Its preimage is closed and contains
$\mathcal N_{a,b}$. Reality specifies a closed real form.
\end{proof}

The quotient removes all carrierwise constants and all other
representations of the zero normalized total. It does not select
a physical vacuum counterterm; a vacuum sector, if needed, stays
separate. Equality of germs on every finite-coordinate subspace
is equivalent to \eqref{aqb-null-space} by analytic uniqueness.

Use the cell norm $\|A\|_1=\sum_i\|A|_i\|_\infty$.
For $\|A\|_\infty,\|B\|_\infty\le\rho<b$, directional
Cauchy on the convex segment gives
\begin{align}
 |E_H(A|_H)-E_H(B|_H)|
 &\le\frac{\|E_H\|_b}{b-\rho}
                          \|(A-B)|_H\|_\infty\notag\\
 &\le\frac{\|E_H\|_b}{b-\rho}
                          \sum_{i\in H}\|A_i-B_i\|_\infty.
 \label{aqb-local-Lipschitz}
\end{align}
Use a strictly smaller directional circle first, then take its
radius to the displayed margin. Summing at each cell proves
\begin{equation}
 |\Phi_E(A)-\Phi_E(B)|
       \le\frac{\|E\|_{a,b}}{b-\rho}\|A-B\|_1.
 \label{aqb-summable-Lipschitz}
\end{equation}
Thus \eqref{aqb-total} extends by an absolutely convergent series
to summable fields in every strict interior ball. Finite field
truncations converge in \eqref{aqb-summable-Lipschitz}. A null
family remains zero on these fields. The extension therefore
depends only on the quotient class, and the estimate holds with
the quotient norm after taking the infimum over representatives.

\subsection{A continuous total Maxwell coefficient}
\label{aqb-beta-section}

For each fixed order $n\le4$, scalar Cauchy and the exact finite
monomial count give a summable total derivative kernel. Indeed
the ordered derivative sum on a carrier is at most
$(D|H|+n)^n b^{-n}\|E_H\|_b$. Connectedness bounds every
slot displacement by $c_dJ|H|$. Spending half the carrier
exponent on spatial distance and half on the fixed polynomial
gives, for sufficiently small fixed $\sigma=\sigma(a,J)>0$,
\begin{equation}
 \sup_{i_1}\sum_{i_2,\ldots,i_n}
 e^{\sigma\sum_{j\ge2}|x_j-x_1|}
 |K_n(E;i_1,\ldots,i_n)|
       \le C_n(a,D,J)b^{-n}\|E\|_{a,b}.
 \label{aqb-kernel-bound}
\end{equation}
All scalar and color slots are included. Fixed moments follow
with a further exponent loss. This is a fixed-order estimate,
not a dimension-free absolute Taylor estimate at all orders.

Each kernel component is a derivative on a finite-coordinate
subspace and hence depends only on $[E]$. Taking a representative
infimum gives \eqref{aqb-kernel-bound} with $\|[E]\|_{\mathcal Q}$.
For a simple gauge group and an invariant total, define in the
trace-normalized color component
\begin{equation}
 \beta([E])=\frac12\partial_{p_\nu}^2
                       \Gamma^E_{2,\mu\mu}(0),
 \qquad \mu\ne\nu.
 \label{aqb-beta-definition}
\end{equation}
Use positive-bond midpoints $\xi(x,\mu)=x+e_\mu/2$.
The total Ward algebra of Theorem~\ref{tkw-theorem} applies
directly to these infinite kernels. It identifies
\eqref{aqb-beta-definition} as the Maxwell coefficient and makes
the distinct pair immaterial. No subsequence is needed once
the infinite interaction is specified. For real interactions
this coefficient is real.

It is a continuous linear functional on the invariant quotient:
\begin{equation}
 |\beta([E])|\le C_\beta(a,D,J)b^{-2}
                                      \|[E]\|_{\mathcal Q}.
 \label{aqb-beta-bound}
\end{equation}
For example, a second derivative component is bounded by
$2b^{-2}\|E_H\|_b$, there are at most $D|H|$ second slots,
and their squared midpoint separation is at most
$(c_dJ|H|)^2$. Thus, up to a fixed color-normalization factor,
one can take
\begin{equation}
 C_\beta=D(c_dJ)^2\sup_{m\ge1}m^3e^{-am}.
 \label{aqb-beta-constant}
\end{equation}
The factor $b^{-2}$ is essential. For a semisimple product use
the finite vector of coefficients and one reference per simple
ideal.

\subsection{Full-reference extraction without a small-radius loss}
\label{aqb-extraction-section}

Let $[S]$ be an invariant reference with $\beta([S])=1$.
Assume a representation on one fixed radius $\rho$ satisfies
$\|S\|_{a,\rho}\le M_S$, uniformly in preceding refinement,
and its total first derivative vanishes. This is the uniform
fixed-radius classical support input; normalization of the
Hessian alone does not imply it.

Define the homogeneous phase projections coefficientwise by
\begin{equation}
 (P_nE_H)(A)=\frac1{2\pi}\int_0^{2\pi}
                     e^{-int}E_H(e^{it}A)\,dt.
 \label{aqb-phase-projection}
\end{equation}
They retain $H$ and have norm at most one. Replace $S_H$ by
$S_H^{\ge2}=\sum_{n\ge2}P_nS_H$. The normalized total has
the same class, since constants are removed and the total
linear sector is zero. Homogeneity yields, for $b\le\rho/2$,
\begin{equation}
 \|S^{\ge2}\|_{a,b}
 \le\frac{(b/\rho)^2}{1-b/\rho}M_S
 \le C_Sb^2,
 \qquad C_S=\frac{2M_S}{\rho^2}.
 \label{aqb-reference-bound}
\end{equation}
No vanishing of individual linear coefficients or individual
gauge invariance was assumed. A joint subsequential limit
preserves this bound by Fatou.

The full-reference projection
\begin{equation}
 \mathcal P_S[E]=[E]-\beta([E])[S]
 \label{aqb-extraction}
\end{equation}
satisfies $\beta\mathcal P_S=0$, $\mathcal P_S^2=\mathcal P_S$,
and
\begin{equation}
 \|\mathcal P_S\|\le C_{\rm ext},
 \qquad C_{\rm ext}=1+C_\beta C_S.
 \label{aqb-extraction-bound}
\end{equation}
The constant is independent of small $b$. Apply
\eqref{aqb-beta-bound} and \eqref{aqb-reference-bound} to
representatives approaching the quotient infimum. The whole
invariant reference is subtracted, so its nonlinear Ward
identities remain exact.

\subsection{Classical substitution on the quotient}
\label{aqb-descent-section}

Fix $L$ and an actual infinite pre-gauge background map $a_L$.
Assume its coherent whole-fiber support family satisfies
\begin{equation}
 a_L|_i=\sum_{X\ne\varnothing}a_{i,X}(w_X),\qquad
 a_{i,X}(0)=0,\qquad
 \sup_i\sum_X e^{\kappa\tau(\{i\}\cup X)+A|X|}
                \|a_{i,X}\|_b\le Cb/L,
 \label{aqb-background-family}
\end{equation}
and $\|a_L(w)\|_\infty\le\rho_0<b$. The maximum over the
whole current output fiber is inside the support sum, and
overlapping collars use one coefficient assignment. For compact
$w$ supported in $K$, a nonzero term has $X\cap K\ne\varnothing$;
hence
\begin{equation}
 \|a_L(w)|_i\|\le(Cb/L)e^{-\kappa\operatorname{dist}(i,K)}.
 \label{aqb-background-summability}
\end{equation}
At fixed $L$ there are finitely many current coordinates per
fiber, so $a_L(w)$ is summable. Consequently
\begin{equation}
 (\mathcal T_L\Phi_E)(w)=\Phi_E(a_L(w))
 \label{aqb-scalar-substitution}
\end{equation}
depends only on $[E]$, by \eqref{aqb-summable-Lipschitz}.
This is algebraic classical descent.

For bounded descent into an activity quotient, suppose a
connected strict representation $\mathscr L_LE$ of
\eqref{aqb-scalar-substitution} obeys
\begin{equation}
 \|\mathscr L_LE\|_{a_{\rm out},b}
                         \le B_L\|E\|_{a,b}.
 \label{aqb-representation-input}
\end{equation}
Then $\mathcal T_L[E]=[\mathscr L_LE]$ is well defined and has
norm at most $B_L$. Distinct valid output assignments have the
same normalized total; equivalent inputs give equal output
totals by summable-field continuity. The class map is linear
because its scalar total is linear in $E$, even without a linear
choice of representatives. Taking near-minimal input
representatives proves the norm bound. Scalar amplitude alone
does not prove \eqref{aqb-representation-input}.

The compactness construction of Section~\ref{aic-section}
provides simultaneous infinite interactions and coherent maps
under its fixed-parameter finite-regulator hypotheses. Its
summable-field composition limit transfers finite classical
identities only on one jointly chosen system, with common strict
domain margins. Unrelated subsequences cannot be composed.
Preserving invariant totals needs more than an abstract orbit
relation. The connecting infinite gauge must be one to which
the input total's gauge identity applies. A sufficient condition
is approximation by allowed compact gauges with transformed
fields converging in $\ell^1$ inside a common strict margin;
then \eqref{aqb-summable-Lipschitz} passes the identity. A joint
finite invariant approximation of a particular interaction and
the maps is another sufficient condition for that interaction.
For an action on the entire invariant quotient, this requirement
must hold for every input class. A joint approximation of one
input, or unrelated finite subsequences, does not establish a
universal invariant-quotient action.

\subsection{Conditional classical contraction after extraction}
\label{aqb-contraction-section}

The following statement keeps the infinite-family and scale
requirements explicit. Fix $J,a$ and the primitive block ratio.
Use one joint infinite system with normalized references $S_M$
and $S_{ML}$ satisfying \eqref{aqb-reference-bound} uniformly.
Assume the coherent physical $C^{1,\alpha}$ families of
Proposition~\ref{pgh-support-proposition}, with primitive
constants independent of $L,M$, and impose the fixed support
reserves of Theorem~\ref{lst-theorem} after choosing $J$.
Assume also the core/tail budgets of the fifth-order background
transport theorem of Notebook28, with unchanged $a,b$.
Its constant $C_{\rm high}$ is independent of $L$ once these
reserves are fixed. Increasing $L$ weakens its compression
requirement; the same fixed $J$ and logarithmic tail cutoff
remain admissible. Choose $b$ after the final finite $L$, so
$b\le r_A/L_*$ with $L_*=L/(8C_A)$ and all other classical
margins hold. Scalar substitution and preservation of the
invariant domain satisfy the additional gauge-identity requirement
just specified. An operator assertion on the whole invariant
quotient assumes it for every class; a joint finite approximation
may instead verify the following pointwise estimate on its
specified collection of classes.

\begin{theorem}[Classical quotient estimate]
\label{aqb-classical-theorem}
Under these hypotheses, in dimension four,
\begin{equation}
 \|\mathcal T_L\mathcal P_{S_M}[E]\|_{\mathcal Q}
 \le C_{{\rm ext,in}}
       \bigl[C_{\rm low}(J)L^{-\alpha}
                    +6C_{\rm high}L^{-1}\bigr]
                     \|[E]\|_{\mathcal Q}.
 \label{aqb-classical-bound}
\end{equation}
If a normalized output reference is used, then
\begin{equation}
 \|\mathcal P_{S_{ML}}\mathcal T_L\mathcal P_{S_M}\|
 \le C_{{\rm ext,out}}C_{{\rm ext,in}}
       \bigl[C_{\rm low}(J)L^{-\alpha}
                    +6C_{\rm high}L^{-1}\bigr].
 \label{aqb-projected-contraction}
\end{equation}
The constants are fixed after $J$ and before $L,b$. A sufficiently
large finite allowed $L=R^q$, followed by a sufficiently small
common positive $b$, makes the last bound less than one.
\end{theorem}

\begin{proof}
Put $[F]=\mathcal P_{S_M}[E]$ and choose any representative.
The phase projections \eqref{aqb-phase-projection} retain carriers,
have norm at most one, and give a coefficientwise fifth-order
remainder $R_5=I-\sum_{n=0}^4P_n$ with norm at most six.
These projections also descend to the unrestricted quotient:
apply the phase integral to a zero normalized analytic total.
They need not preserve gauge invariance separately.

The total linear term of $F$ is zero by global adjoint invariance,
its vacuum is removed, and $\beta(F)=0$. The total Ward algebra
therefore gives the low-kernel cancellations in
Theorem~\ref{lst-theorem}. The kernel bound
\eqref{aqb-kernel-bound} supplies that theorem's $\epsilon$ with
at most a fixed $C_{\rm ker}(J)\|F\|_{a,b}$. Its new connected
representation has norm at most
$C_{\rm low}(J)L^{-\alpha}\|F\|_{a,b}$.
The coefficientwise $R_5$ part has a connected representation
of norm at most $6C_{\rm high}L^{-1}\|F\|_{a,b}$ by the
fifth-order transport theorem. The zero total linear term stays
zero after substitution by \eqref{aqb-summable-Lipschitz}; it
need not vanish on each carrier.

All scalar decompositions are legitimate on summable inputs:
the phase projections and $R_5$ lie in the anchored space,
Section~\ref{aqb-quotient-section} gives their continuous
extensions, and the new output families converge normally on
compact fields. Adding the two representations and taking the
infimum over representatives of $[F]$ gives the bracket in
\eqref{aqb-classical-bound}. Apply
\eqref{aqb-extraction-bound} at the input, and if desired at
the output, to obtain both displays.
\end{proof}

The output projection is explicit: the theorem does not claim
that classical substitution itself preserves $\beta=0$.
The scalar semigroup of Theorem~\ref{csg-theorem}, passed on
a joint system as in Section~\ref{aic-section}, permits the
finite choice $L=R^q$ after $R,J$ are fixed. The bound is on a
class admitting a newly justified localization, not on an
arbitrarily fixed sequential carrier assignment.

\subsection{Quantum descent is a separate obligation}
\label{aqb-scope}

No quantum map on $\mathcal Q_{a,b}$ has been defined here.
Equality of infinite normalized totals on compact fields does
not by itself identify finite-regulator exponentials, their
sharp restricted integrals, chosen vacuum normalizations, or
their limits. Thus this quotient does not yet permit arbitrary
equivalent representatives to replace the actual quantum output
inside its next finite integration. Quantum descent and its
continuity and normalization compatibility, the full amplitude
recursion, and the coupling trajectory remain separate.
